\section*{Limitations}

\paragraph{Synthetic construction.}
\benchmark{} is synthetic: it combines human-authored business scenarios, document schemas, accounting policies, tax and FX assumptions, support-document dependencies, distractors, and contradiction templates with programmatic sampling, rendering, and ledger replay. This design provides inspectable labels and controlled counterfactuals but limits document fidelity. Real accounting document bundles contain messier layouts, OCR errors, missing pages, handwriting, ambiguous vendor formats, inconsistent naming, and institution-specific artifacts that our rendered documents only approximate. The benchmark measures reasoning over templated OCR text under these controlled conditions. The character-corruption control covers seeded text noise and excludes raw OCR-engine output and rendered layout. Consequently, the experiments leave the contribution of flattened layout to document-linking failure unresolved.

\paragraph{Accounting coverage.}
Accounting is broad and policy-dependent. Our labels are deterministic after fixing the chart of accounts, accounting policy, tax assumptions, and transaction generator. We score one canonical treatment under the visible policies. Equivalent decompositions, presentation variants, and other professionally acceptable treatments remain outside the current exact-match scope. We cover multi-document postings, period-end adjustments, tax, foreign exchange, leases, deferred tax, revenue schedules, subledgers, distractors, and forced contradictions. The current scope excludes the full range of GAAP, IFRS, local tax regimes, materiality judgments, industry-specific policy choices, and audit procedures. Future policy variants could enumerate multiple acceptable-answer sets. The benchmark defines a precise document-grounded accounting setting within this fixed policy space.

\paragraph{Analysis depth.}
The ablations show several patterns, including the aggregation gap, weak gains from prompt-only changes, document-linking failures, context/evidence effects, and the ledger-feedback trade-off with inconsistency detection. The analysis is behavioral and leaves causal mechanisms and cross-domain transfer open. The released evaluation framework supports follow-up studies with broader ablations, controlled perturbations, model-family comparisons, and mechanistic or representation-level analyses of why models fail after the relevant evidence is present.

\paragraph{Evaluation scope.}
Headline numbers use a single temperature-$0$ decoding run except for the best-of-3 condition, so the paper emphasizes large paired effects over single-point differences. Six models are evaluated on the full 710-record core split. Targeted ablations use the compact 143-record split, which preserves coverage over every industry--period--difficulty cell and inconsistency code. Repeated-run controls measure stability, while commercial APIs may drift despite pinned request/response identifiers. Released outputs support exact offline rescoring, whereas future API calls may differ. The evidence supports an aggregation gap for this evaluated model panel. Broader generalization requires additional model families.

\paragraph{Benchmark longevity.}
The public core split may become contaminated as released benchmarks enter model training data. \benchmark{} therefore includes a generator alongside the fixed test file. Users can sample new seeded document bundles from the same human-authored scenario space, create larger train/test splits, and add new industries, document schemas, policies, contradiction templates, and ablation conditions without hand-labeling each record. A fixed release enables comparison across studies, while periodically generated held-out bundles and community extensions can reduce contamination over time.

\paragraph{Intended use and risk.}
\benchmark{} evaluates research systems under synthetic conditions and does not certify accounting competence. Professional review remains necessary when selecting or deploying systems for bookkeeping, audit, tax, regulatory reporting, or financial decision-making. Applying model outputs to real records without qualified human oversight could propagate incorrect postings, overlook contradictory evidence, or produce materially misleading statements. The synthetic scope and exact-match metrics leave those deployment risks unmeasured.

\paragraph{Ethics and AI assistance.}
The benchmark uses synthetic data generated without real records, PII, or production accounting data. \appref{appx:reproducibility} reports reproducibility and release details. Generative AI was used for coding assistance and language, wording, and formatting edits. The authors developed and verified the benchmark design, accounting logic, experiments, analyses, interpretations, and final claims.
